\begin{proposition}[Every secondary Hopf surface of algebraic dimension zero occurs unchanged]\label{cdp:st:C.2.UN.43}
Every secondary Hopf surface $H$ of algebraic dimension zero
has an actual finite quotient presentation of one of the forms
\[
\begin{array}{ll}
 H_{\alpha,\beta}/\langle(z,w)\mapsto
          (\zeta_n^a z,\zeta_n^b w)\rangle,
 &\gcd(a,n)=\gcd(b,n)=1,\\[2pt]
 H_r(\beta)/\langle(z,w)\mapsto
          (\zeta_n^r z,\zeta_n w)\rangle,
 &\gcd(r,n)=1.
\end{array}
\tag{C.SH.1}\label{cdp:eq:C.SH.1}
\]
The first primary contraction is $(\alpha z,\beta w)$ with no
multiplicative resonance, and the second is
$(\beta^r z+w^r,\beta w)$, where $0<|\alpha|,|\beta|<1$.
This is an exhaustive presentation, without a uniqueness
assertion about the displayed parameters.

Every prescribed such surface occurs itself as a smooth
multiplicity-one component of a reduced SNC fibre of a proper
map from a smooth compact complex threefold to a smooth
compact curve, with connected fibres and general fibre a
complex two-torus. All reduced fibre supports can have Euler
number zero. In the diagonal construction the base is the
finite torsion-cover curve $B'$ constructed below; no claim
that this curve is rational is needed.
\end{proposition}
\Uses{the hypotheses, notation and definitions stated here; \cite[Theorem~2.1]{Khaled2024Hopf}; \cref{cdp:st:C.2.UN.40}; \cref{cdp:eq:C.SH.3}; \cite[\S3, pp.~230--232]{Kato1975Hopf}; \cref{cdp:st:C.2.UN.41}; \cref{cdp:eq:C.SH.9}; \cref{cdp:eq:C.SH.10}; \cref{cdp:st:C.2.UN.1}; \cref{cdp:st:C.2.UN.4}; \cref{cdp:st:C.2.UN.42}}


\begin{proof}
\par\medskip\noindent\textbf{The actual deck group.}
Write $W=\mathbb C^2\setminus\{0\}$ and $H=W/\Gamma$.
The deck group has an infinite cyclic contraction subgroup of
finite index. Its normal core, the intersection of its finitely
many conjugates, is a normal infinite cyclic subgroup
$\langle F\rangle$ of finite index. Conjugation acts on it by
$1$ or $-1$. The nonzero conjugation-invariant character
$\gamma\mapsto\log|\det(d\gamma)_0|$ excludes the second
possibility. Thus $F$ is central and the actual primary cover is
\[
                 P=W/\langle F\rangle\longrightarrow H.
\tag{C.SH.2}\label{cdp:eq:C.SH.2}
\]
Taking elementary symmetric functions on its finite generic
sheets gives $a(P)=a(H)=0$.

The primary normal forms and their full centralizers are
\cite[Theorem~2.1]{Khaled2024Hopf}. In the nonresonant diagonal
case a commuting power series has the form $(Az,Bw)$:
the eigenmonomial argument of \ref{cdp:st:C.2.UN.40} leaves only $z$ and
$w$ in its respective coordinate functions. In the
non-diagonal case the full centralizer is
\[
       (z,w)\longmapsto(A^r z+B w^r,Aw),\qquad A\ne0.
\tag{C.SH.3}\label{cdp:eq:C.SH.3}
\]
Indeed write a commuting map as $(P,Q)$. Its second
coordinate satisfies $Q\circ F=\beta Q$. The weighted
eigenfunction calculation of \ref{cdp:st:C.2.UN.40} forces $Q=Aw$.
The first coordinate then satisfies
$P\circ F=\beta^rP+A^r w^r$. On a weighted-degree-$d$
piece with $d\ne r$, the operator $F^*-\beta^r$ is
invertible, since its only eigenvalue is
$\beta^d-\beta^r\ne0$. Thus $P=Cz+Bw^r$.
Comparing its $w^r$ coefficients gives $C=A^r$;
invertibility requires $A\ne0$. This proves the full
formula \eqref{cdp:eq:C.SH.3} directly.
The parameters $(A,B/A^r)$ in this formula have product
$(A_1,c_1)(A_2,c_2)=(A_1A_2,c_1+c_2)$. Both centralizers
are abelian. Since every element of $\Gamma$ commutes
with $F$, the actual deck group is abelian.

A finitely generated abelian group containing $\mathbb Z$
with finite index has the form $\mathbb Z\oplus T$, with $T$
its finite torsion subgroup. Choose a generator $f$ of the
infinite factor with the contracting sign. Some positive power
is $F^k$ times an element of $T$, $k>0$. Thus both diagonal
eigenvalue moduli of $f$ are less than one. In \eqref{cdp:eq:C.SH.3} it has
$|A|<1$ and nonzero additive parameter, since the parameter
of $F$ is nonzero. Rescaling $z$ makes its resonant coefficient
one. We obtain the actual holomorphic quotient maps
\[
       H=(W/\langle f\rangle)/T,\qquad
       \Gamma=\langle f\rangle\times T.
\tag{C.SH.4}\label{cdp:eq:C.SH.4}
\]
For diagonal $T\subset(\mathbb C^*)^2$, each coordinate
character is injective: its kernel would fix a nonzero point
of the corresponding axis. Hence $T$ is cyclic and both
weights are units modulo its order. In \eqref{cdp:eq:C.SH.3}, a finite-order
element has additive parameter zero and acts as
$(A^r z,Aw)$. Its finite subgroup is cyclic, and freeness
on the first axis forces $\gcd(r,n)=1$. This proves \eqref{cdp:eq:C.SH.1}
and the finite-primary-cover comparison, including the
exclusion of a nonabelian group in algebraic dimension zero.
The abelian calculation also appears in
\cite[\S3, pp.~230--232]{Kato1975Hopf}; no topological
classification is used here.

Conversely the indicated finite maps are unramified.
An equation $\tau^k(z,w)=f^\ell(z,w)$ at a nonzero coordinate
forces $\ell=0$ by the modulus of the contraction multiplier;
the coprimality conditions then force $k=0$ modulo $n$.
In the non-diagonal case use $w$ when it is nonzero and
$z$ on its complement. Thus these are the exact free
finite actions to be extended.

\par\medskip\noindent\textbf{A prescribed torsion section at the nodal fibre.}
For the diagonal case start with the compact elliptic surface
$J\to\mathbb P^1$, zero section, nodal fibre at zero and
section $P$ of \ref{cdp:st:C.2.UN.41}. Its multiplicative value at that
fibre is the prescribed $\alpha$. Explicitly, on
$y^2=x^2(x+1)$ the point with group coordinate $\alpha$ is
\[
 \left(\frac{4\alpha}{(1-\alpha)^2},
       \frac{4\alpha(1+\alpha)}{(1-\alpha)^3}\right).
\tag{C.SH.5}\label{cdp:eq:C.SH.5}
\]
Blowing up the nine distinct smooth base points of a general
cubic pencil through this point and the origin gives that
surface and its section, with an unchanged $I_1$ fibre.

Let $B^\circ$ be the smooth elliptic-fibre locus.
The exact-order-$n$ part of $J[n]|_{B^\circ}$ is a finite
etale algebraic cover. Select a component by the following
actual local section. The smooth group locus of the nodal
family has central group $\mathbb C^*$; multiplication by
$n$ has nonzero derivative at $\zeta_n$. The holomorphic
implicit-function theorem gives a unique local $n$-torsion
section specializing to $\zeta_n$, necessarily of exact order.
Take the connected cover containing its punctured restriction,
and normalize $\mathbb P^1$ in its function field. This gives
a smooth connected projective curve and finite map
\[
                             B'\longrightarrow\mathbb P^1.
\tag{C.SH.6}\label{cdp:eq:C.SH.6}
\]
At the selected point $b_0$ this map is unramified: the local
sheet is the graph of the just constructed section and has
degree one over the disc. The actual $I_1$ fibre, its simple
smoothing parameter and $P(b_0)=\alpha$ remain unchanged.

Resolve the compact algebraic base-changed elliptic surface
and contract vertical $(-1)$-curves, away from this unchanged
tube, to obtain a smooth relatively minimal surface $J'\to B'$.
The pulled-back section and the tautological torsion point
extend as sections $P',R$: rational maps from the smooth
projective base curve to the proper regular surface extend by
the valuative criterion. They satisfy
\[
                 P'(b_0)=\alpha,\qquad R(b_0)=\zeta_n,
                 \qquad \operatorname{ord}(R)=n.
\tag{C.SH.7}\label{cdp:eq:C.SH.7}
\]
Their translations extend to automorphisms of $J'$.
The minimal-graph argument of \ref{cdp:st:C.2.UN.41} applies because
$K_{J'}\cdot C\ge0$ for every vertical curve, by adjunction
and relative minimality, or by Kodaira's fibre components.
A $(-1)$-curve exceptional for one graph projection but not
the other would have nonnegative intersection with its
canonical pullback and all other effective discrepancy
curves, contradicting canonical intersection $-1$.
Both translations and inverses are therefore holomorphic;
their commutation on the smooth fibres proves commutation
everywhere.

\par\medskip\noindent\textbf{The diagonal action and the actual compact quotient.}
Put $q_0=\alpha\beta$ and form the compact suspension
\[
 Y^{\rm std}=(\mathbb C^*_u\times J')/
       \langle(u,y)\mapsto(q_0u,y+P')\rangle.
\tag{C.SH.8}\label{cdp:eq:C.SH.8}
\]
Its proper compact bundle comparison is the one in
\ref{cdp:st:C.2.UN.41}. The commuting order-$n$ automorphism
\[
               (u,y)\longmapsto
                  (\zeta_n^{a+b}u,y+aR)
\tag{C.SH.9}\label{cdp:eq:C.SH.9}
\]
descends. On the selected Tate tube $R=[\zeta_n]$ exactly,
by uniqueness of the torsion section in the smooth group
locus. In the multiplicative coordinate change $z=y$, $w=u/y$,
\eqref{cdp:eq:C.SH.9} is precisely
\[
                    (z,w)\longmapsto
                    (\zeta_n^a z,\zeta_n^b w).
\tag{C.SH.10}\label{cdp:eq:C.SH.10}
\]
Replace the tube by \ref{cdp:st:C.2.UN.41}'s two-Hopf plumbing using
its exact punctured-family isomorphism. The action extends
on the first Hopf copy by \eqref{cdp:eq:C.SH.10}, on the second by the
inverse weights, and on the first neck by
\[
 (z,x,y)\longmapsto
     (\zeta_n^a z,\zeta_n^b x,\zeta_n^{-b}y),
                         \qquad xy=t.
\tag{C.SH.11}\label{cdp:eq:C.SH.11}
\]
Interchange $a,b$ and the axes for the other neck.
Substitution in $(z,w)\mapsto(z^{-1},t/w)$ and
$(z,w)\mapsto(t/z,w^{-1})$ proves equivariance of both
gluing maps. The action is free on the whole central
support: on a double curve the equation
$\zeta_n^{ak}z=\alpha^\ell z$, or its counterpart, forces
$\ell=0$ by modulus and then $k=0$ by coprimality.
Off the double curves use the already verified Hopf
freeness. On a nearby torus it is a finite translation,
also free. Hence the quotient is smooth on this entire
tube and its two central components are the prescribed
secondary Hopf surface itself. The inverse weights
generate the same subgroup.

Take the finite quotient of the entire compact threefold
and resolve it over the other finitely many bad values,
leaving this smooth tube unchanged. General fibres are
quotients of tori by translations, hence tori. Connectedness
passes through the finite quotient and the connected fibres
of the birational resolution. The distinguished reduced
SNC fibre has two Hopf components and two elliptic double
curves, and therefore Euler number zero. Every other
reduced fibre has even Euler number by \ref{cdp:st:C.2.UN.1}.
Apply \ref{cdp:st:C.2.UN.4} only over those other values to make their
Euler numbers zero; none of these operations changes
the specified Hopf surface or its distinguished tube.
This proves compact diagonal secondary occurrence over
the actual curve $B'$.

\par\medskip\noindent\textbf{The nonlinear scalar action.}
Use the intermediate compact family $\widetilde Y\to C$ in
\eqref{cdp:eq:C.NH.11} for the prescribed $H_r(\beta)$, before the
additional involution in \eqref{cdp:eq:C.NH.15}. Its smooth nonzero
tautological line space has coordinates
\[
 (z,w,x)\in(\mathbb C^2\setminus\{0\})\times\mathbb C,
                   \qquad p-p_0=x^r z,\quad s=xw.
\tag{C.SH.12}\label{cdp:eq:C.SH.12}
\]
Multiplication of that line by $\zeta_n$ is the actual map
\[
 (z,w,x)\longmapsto
       (\zeta_n^r z,\zeta_n w,\zeta_n^{-1}x).
\tag{C.SH.13}\label{cdp:eq:C.SH.13}
\]
It commutes with the lifted base translations and the
contracting deck map. It is free when $\gcd(r,n)=1$:
a nonzero $w$ or $x$ forces the corresponding root to be
one; if both vanish then $z\ne0$ and its $r$-th power
must be one. It stays free after the infinite cyclic
quotient. Away from the exceptional component the
canonical line trivialization turns equality with a
deck power into $\zeta_n^k=\beta^\ell$, so $\ell=0$.
On the Hopf component use its nonzero $w$ coordinate,
or its $z$ coordinate if $w=0$, for the same modulus
argument. Coprimality then finishes the proof.

The resulting quotient is a smooth compact threefold.
Its desired component is exactly the second quotient
in \eqref{cdp:eq:C.SH.1}, with no change to that surface. On every
general torus fibre the action is a line-coordinate
translation, so the quotient is a torus. The distinguished
full fibre remains reduced SNC, and its Euler number
zero is preserved by the free finite quotient. Any
remaining values can be balanced by \ref{cdp:st:C.2.UN.4} away
from that fibre.

The base for this construction is the specified compact
curve in \ref{cdp:st:C.2.UN.42}. No rational-base globalization is
needed for this occurrence assertion.
\end{proof}